% 同一插值直线上的三个参数；从零初始化的梯度下降停在欧氏投影。
% 轨迹画在w_1=w_2上；细弧是范数1/sqrt(2)的等值线。
\begin{tikzpicture}[figure picture,foundation picture]
 \draw[foundation axis] (12,8)--(65,8) node[right] {$w_1$};
 \draw[foundation axis] (12,8)--(12,62) node[above] {$w_2$};
 \node[below left] at (12,8) {$0$};
 \draw[draw=FigureMuted,line width=.6pt] (43.82,8) arc[start angle=0,end angle=90,radius=31.82];
 \draw[FigureBlue,line width=1.1pt] (12,53)--(57,8);
 \node[anchor=west] at (38,47) {$w_1+w_2=1$};
 \draw[FigureRed,line width=1.1pt,-{Stealth[length=1.8mm]}] (12,8)--(34.5,30.5);
 \foreach \pos in {23.25,28.875,31.6875}{
  \fill[FigureRed] (\pos,\pos-4) circle (.75);
 }
 \node[align=center,text=FigureRed] at (33,17) {$\bm w^{(0)}=\bm0$\\$w_1=w_2$};
 \fill[FigureBlue] (57,8) circle (1.1);
 \fill[FigureYellow] (34.5,30.5) circle (1.5);
 \fill[FigureBlue] (12,53) circle (1.1);
 \node[below] at (57,7) {A};
 \node[anchor=west] at (37,30.5) {B};
 \node[left] at (11,53) {C};

 \node[anchor=west] at (72,43) {A：$(1,0)^\top\mapsto1$};
 \node[anchor=west] at (72,28) {B：$(\tfrac12,\tfrac12)^\top\mapsto\tfrac12$};
 \node[anchor=west] at (72,13) {C：$(0,1)^\top\mapsto0$};
\end{tikzpicture}%
